\unitopen{32}
\begin{candos}
\item wyrazić opinię, zgodę i sprzeciw uprzejmie, formalnie i potocznie;
\item ułożyć krótki argument z przykładem (\de{erstens, zweitens, außerdem, zum Beispiel});
\item użyć łączników \de{deshalb, deswegen, trotzdem, allerdings, jedoch} z poprawnym szykiem;
\item uczestniczyć w dyskusji na spotkaniu technicznym.
\end{candos}

\section{Słownictwo}
\vocabtable{32}

\section{Wyrażanie opinii}
Opinia: \de{Ich finde, dass \dots} \de{Ich denke, \dots} \de{Meiner Meinung nach \dots} \de{Ich bin der Meinung, dass \dots}
Zgoda: \de{Da stimme ich zu.} \de{Genau.} \de{Das sehe ich auch so.} Częściowa zgoda: \de{Das stimmt zwar, aber \dots}
Sprzeciw: \de{Da bin ich anderer Meinung.} \de{Das sehe ich anders.} \de{Ich bin nicht sicher, ob das stimmt.}

\section{Łączniki i szyk}
\emph{Przysłówki łączące} (\de{deshalb, deswegen, darum, trotzdem, außerdem, allerdings, jedoch}) \textbf{nie zmieniają} szyku
zdania podrzędnego, ale zajmują \emph{pierwsze miejsce}, więc po nich stoi \textbf{czasownik}:
\de{Der Server ist langsam, \textbf{deshalb haben} wir einen Cache eingeführt.} \de{Das ist teuer, \textbf{trotzdem kaufen} wir es.}
Spójniki współrzędne \de{und, aber, oder, denn, sondern} stoją przed zdaniem bez wpływu na szyk.

\begin{zasada}[Struktura argumentu]
(1) stanowisko: \de{Ich finde, dass wir den Cache einführen sollten.} (2) uzasadnienie: \de{Erstens ist die Latenz zu hoch.
Zweitens sparen wir Kosten.} (3) przykład: \de{Zum Beispiel dauert die Abfrage sechs Sekunden.} (4) zastrzeżenie:
\de{Allerdings müssen wir die Konsistenz prüfen.} (5) wniosek: \de{Deshalb schlage ich vor, dass wir einen Test machen.}
\end{zasada}

\begin{falle}
Polskie „dlatego” nie zmienia szyku, więc łatwo napisać \de{deshalb wir haben} (błąd). Poprawnie: \de{deshalb haben wir}.
Drugi błąd: \de{weil} i \de{deshalb} to nie to samo: \de{weil} wprowadza zdanie podrzędne (czasownik na końcu), \de{deshalb}
jest przysłówkiem (czasownik drugi). Trzeci: \de{sondern} (po przeczeniu) vs \de{aber}: \de{Nicht heute, sondern morgen.}
\end{falle}

\begin{szwajcaria}
W szwajcarskiej kulturze dyskusji ceni się \emph{konsensus} i umiar: bezpośrednie \de{Das ist falsch} bywa odbierane jako
agresywne. Łagodniej: \de{Ich sehe das etwas anders} lub \de{Ich hätte da eine andere Idee.} Spotkania kończą się często
\de{Wir sollten das nochmals anschauen}, co oznacza „wrócimy do tego”, a nie decyzję.
\end{szwajcaria}

\begin{praca}
\de{Ich schlage vor, dass wir die Datenbank wechseln. Erstens ist die aktuelle zu langsam. Zweitens sind die Lizenzkosten hoch.
Allerdings müssen wir die Migration planen.} \de{Da stimme ich dir zu, aber wir sollten den Aufwand schätzen.}
\end{praca}

\section{Dialogi}
\textbf{Dialog 1 --- dyskusja o technologii}
\begin{dialog}
\turn{A}{Ich finde, dass wir auf Azure wechseln sollten. Meiner Meinung nach ist es günstiger.}{Uważam, że powinniśmy przejść na Azure. Moim zdaniem jest taniej.}
\turn{B}{Das stimmt zwar, allerdings haben wir noch keine Erfahrung damit.}{To prawda, jednak nie mamy z tym jeszcze doświadczenia.}
\turn{A}{Deshalb schlage ich einen Pilot vor. Zum Beispiel könnten wir ein kleines Projekt migrieren.}{Dlatego proponuję pilota. Na przykład moglibyśmy zmigrować mały projekt.}
\turn{B}{Da stimme ich dir zu. Das ist ein guter Kompromiss.}{Zgadzam się. To dobry kompromis.}
\end{dialog}

\section{Ćwiczenia}
\begin{cwiczenia}
\exitem{Wyraź opinię: \de{Ich \dots{}, dass wir testen sollten.} (uważam)}{\de{Ich finde, dass wir testen sollten.}}
\exitem{Wyraź zgodę po niemiecku (formalnie).}{\de{Da stimme ich Ihnen zu.}}
\exitem{Wyraź sprzeciw uprzejmie.}{\de{Da bin ich anderer Meinung.} lub \de{Das sehe ich etwas anders.}}
\exitem{Połącz przysłówkiem: \de{Der Server ist langsam.} (deshalb) \de{Wir führen einen Cache ein.}}{\de{Der Server ist langsam, deshalb führen wir einen Cache ein.}}
\exitem{Połącz przysłówkiem: \de{Es ist teuer.} (trotzdem) \de{Wir kaufen es.}}{\de{Es ist teuer, trotzdem kaufen wir es.}}
\exitem{Popraw: \de{Deshalb wir haben das gemacht.}}{\de{Deshalb haben wir das gemacht.}}
\exitem{\de{aber} czy \de{sondern}: \de{Nicht heute, \dots{} morgen.}}{\de{sondern}}
\exitem{Wskaż: który spójnik wymusza czasownik na końcu: \de{weil} czy \de{deshalb}?}{\de{weil}; po \de{deshalb} czasownik stoi na drugim miejscu}
\exitem{Ułóż wypowiedź: dwa argumenty + zastrzeżenie dla zmiany bazy danych.}{np.\ \de{Erstens ist die Datenbank zu langsam. Zweitens ist sie zu teuer. Allerdings ist die Migration aufwendig.}}
\exitem{Przetłumacz: „Moim zdaniem to jest dobry kompromis.”}{\de{Meiner Meinung nach ist das ein guter Kompromiss.}}
\exitem{Przetłumacz: „Z jednej strony szybciej, z drugiej drożej.”}{\de{Einerseits ist es schneller, andererseits ist es teurer.}}
\end{cwiczenia}

\section{Trening słownictwa (generowany)}
\drill{32}

\weekend{Wybierz kontrowersyjny temat techniczny (np.\ monolit kontra mikroserwisy). Napisz dwuminutowy argument (stanowisko,
dwa uzasadnienia, przykład, zastrzeżenie, wniosek). Powiedz go na głos i poproś AI o kontrargument po niemiecku.}
